\documentclass[12pt,a4paper]{article}
\usepackage[utf8]{vietnam}
\usepackage{amsmath,amssymb}
\usepackage{enumerate}
\usepackage{tabularx}
\usepackage{tikz}
\usepackage{xcolor}
\usepackage[top=0.45cm, bottom=0.45cm, left=0.8cm, right=0.8cm]{geometry}
\DeclareUnicodeCharacter{00B2}{\ensuremath{^2}}
\DeclareUnicodeCharacter{00B3}{\ensuremath{^3}}

\setlength{\parskip}{0pt}
\setlength{\parindent}{0pt}
\setlength{\fboxsep}{3.5pt}
\setlength{\fboxrule}{0.5pt}

% Dòng kẻ chấm mờ cự ly chuẩn viết tay tiểu học (18pt)
\newcommand{\dongke}[1]{%
  \begingroup%
  \setlength{\parskip}{0pt}%
  \baselineskip=18pt%
  \vspace{0.4mm}%
  \foreach \i in {1,...,#1}{%
    \noindent{\color{gray!55}\dotfill}\par%
  }%
  \vspace{0.2mm}%
  \endgroup%
}

% Khung bao quanh bài tập tự luận
\newcommand{\khungBaiTap}[1]{%
  \vspace{0.8mm}%
  \noindent\fbox{%
    \parbox{\dimexpr\linewidth-2\fboxsep-2\fboxrule\relax}{%
      \vspace{0.3mm}%
      #1%
      \vspace{0.3mm}%
    }%
  }%
  \vspace{0.5mm}%
}

% Ô bài tập đặt tính
\newcommand{\oDatTinh}[3][4]{%
  \vspace{0.3mm}%
  \noindent\textbf{#2} #3\par%
  \dongke{#1}%
}

% Câu trắc nghiệm dàn 4 đáp án trên 1 hàng ngang
\newcommand{\cauTN}[6]{%
  \noindent\textbf{Câu #1.} #2\par\vspace{0.1mm}%
  \noindent%
  \makebox[0.25\linewidth][l]{\textbf{A.} #3}%
  \makebox[0.25\linewidth][l]{\textbf{B.} #4}%
  \makebox[0.25\linewidth][l]{\textbf{C.} #5}%
  \makebox[0.25\linewidth][l]{\textbf{D.} #6}\par\vspace{0.6mm}%
}

% Câu trắc nghiệm dàn 2 đáp án / hàng
\newcommand{\cauTNhaiDong}[6]{%
  \noindent\textbf{Câu #1.} #2\par\vspace{0.1mm}%
  \noindent%
  \makebox[0.5\linewidth][l]{\textbf{A.} #3}%
  \makebox[0.5\linewidth][l]{\textbf{B.} #4}\par\vspace{0.1mm}%
  \noindent%
  \makebox[0.5\linewidth][l]{\textbf{C.} #5}%
  \makebox[0.5\linewidth][l]{\textbf{D.} #6}\par\vspace{0.6mm}%
}

\begin{document}
\pagestyle{empty}
\fontsize{10.2pt}{13.2pt}\selectfont
% =============================================================================
% ========================= MÃ ĐỀ: 501 ========================================
% =============================================================================

% -------------------- TRANG 1: MẶT TRƯỚC --------------------
\noindent
\begin{tabularx}{\linewidth}{|X|c|}
\hline
\textbf{PHIẾU BÀI TẬP CUỐI TUẦN -- TOÁN LỚP 5 (TUẦN 16)} & \textbf{MÃ ĐỀ: 501} \\
\textit{Trung tâm Toán Thầy Lĩnh -- Năm học: 2026 - 2027} & \textit{Thời gian: 45 phút} \\
\hline
\end{tabularx}

\vspace{1mm}

\noindent
\begin{tabularx}{\linewidth}{|X|p{3.5cm}|}
\hline
\rule[-2mm]{0pt}{7.5mm}\textbf{Họ và tên:} \dotfill & \textbf{Lớp:} \makebox[2cm]{\dotfill} \\
\hline
\end{tabularx}

\vspace{1mm}

\noindent
\begin{tabularx}{\linewidth}{|p{3cm}|X|}
\hline
\centering\textbf{ĐIỂM SỐ} & \centering\textbf{LỜI NHẬN XÉT CỦA THẦY CÔ} \tabularnewline
\hline
\rule{0pt}{1.1cm} & \tabularnewline
\hline
\end{tabularx}

\vspace{1mm}

\noindent\textbf{BẢNG ĐIỀN ĐÁP ÁN TRẮC NGHIỆM (Mỗi câu đúng 0,25 điểm):}
\vspace{0.6mm}

\begin{tabularx}{\linewidth}{|c|*{12}{>{\centering\arraybackslash}X|}}
\hline
\textbf{Câu} & \textbf{1} & \textbf{2} & \textbf{3} & \textbf{4} & \textbf{5} & \textbf{6} & \textbf{7} & \textbf{8} & \textbf{9} & \textbf{10} & \textbf{11} & \textbf{12} \\
\hline
\textbf{Đáp án} & & & & & & & & & & & & \\
\hline
\end{tabularx}

\vspace{1.2mm}

\noindent\textbf{\large PHẦN I. TRẮC NGHIỆM KHÁCH QUAN (3,0 điểm)}

\cauTN{1}{Kết quả của phép tính $24{,}5 \times 0{,}4$ là:}{$9{,}8$}{$98$}{$0{,}98$}{$9{,}08$}
\cauTN{2}{Giá trị của biểu thức $15{,}6 + 4{,}4 \times 2$ là:}{$40$}{$24{,}4$}{$25{,}4$}{$38{,}8$}
\cauTN{3}{Tỉ số của hai số $3$ và $5$ viết dưới dạng phân số là:}{$\dfrac{5}{3}$}{$3{,}5$}{$\dfrac{3}{5}$}{$0{,}35$}
\cauTN{4}{Phân số $\dfrac{3}{4}$ viết dưới dạng tỉ số phần trăm là:}{$75\%$}{$60\%$}{$80\%$}{$70\%$}
\cauTN{5}{Tỉ số phần trăm của $15$ và $60$ là:}{$20\%$}{$25\%$}{$30\%$}{$15\%$}
\cauTN{6}{$35\%$ của $120\,\text{kg}$ là:}{$42\,\text{kg}$}{$40\,\text{kg}$}{$45\,\text{kg}$}{$35\,\text{kg}$}
\cauTN{7}{Biết $25\%$ của một số là $40$. Số đó là:}{$10$}{$80$}{$160$}{$100$}
\cauTN{8}{Một cửa hàng nhập áo giá $200\,000$ đồng, bán lãi $20\%$. Số tiền lãi là:}{$40\,000$ đồng}{$50\,000$ đồng}{$30\,000$ đồng}{$20\,000$ đồng}
\cauTN{9}{Một hình thang có hai đáy $18\,\text{cm}$ và $12\,\text{cm}$, chiều cao $9\,\text{cm}$. Diện tích hình thang là:}{$270\,\text{cm}^2$}{$135\,\text{cm}^2$}{$145\,\text{cm}^2$}{$150\,\text{cm}^2$}
\cauTN{10}{Một hình tròn có đường kính $10\,\text{cm}$. Diện tích hình tròn đó là:}{$314\,\text{cm}^2$}{$157\,\text{cm}^2$}{$78{,}5\,\text{cm}^2$}{$31{,}4\,\text{cm}^2$}
\cauTNhaiDong{11}{Mảnh đất $500\,\text{m}^2$, dành $40\%$ diện tích làm nhà, phần còn lại trồng cây. Diện tích đất trồng cây là:}{$300\,\text{m}^2$}{$200\,\text{m}^2$}{$250\,\text{m}^2$}{$350\,\text{m}^2$}
\cauTNhaiDong{12}{Lớp 5A có $40$ học sinh, trong đó có $24$ bạn nữ. Tỉ số phần trăm của số bạn nữ so với cả lớp là:}{$55\%$}{$60\%$}{$65\%$}{$50\%$}

\vspace{0.4mm}
\noindent\textbf{\large PHẦN II. TỰ LUẬN (7,0 điểm)}

\noindent\textbf{Bài 1 (2,0 điểm). Đặt tính rồi tính:}
\vspace{0.4mm}

\noindent
\begin{tabularx}{\linewidth}{|X|X|X|X|}
\hline
\oDatTinh[4]{a)}{$48{,}75 + 36{,}8$} & 
\oDatTinh[4]{b)}{$95{,}2 - 47{,}85$} & 
\oDatTinh[4]{c)}{$34{,}5 \times 2{,}6$} & 
\oDatTinh[4]{d)}{$117{,}6 : 4{,}2$} \tabularnewline
\hline
\end{tabularx}

\newpage
% -------------------- TRANG 2: MẶT SAU --------------------
\noindent
\begin{tabularx}{\linewidth}{|X|c|}
\hline
\textbf{PHẦN TỰ LUẬN (TIẾP THEO) -- PHIẾU BÀI TẬP TUẦN 16} & \textbf{MÃ ĐỀ: 501} \\
\hline
\end{tabularx}

\vspace{0.8mm}

% BÀI 2: TÌM X (4 Ý)
\khungBaiTap{%
  \textbf{Bài 2 (1,5 điểm). Tìm $x$:}\par\vspace{0.6mm}
  \noindent
  \begin{minipage}[t]{0.485\linewidth}
    \textbf{a)} $x \times 4{,}5 + 12{,}5 = 39{,}5$\par
    \dongke{3}
    \vspace{0.4mm}
    \textbf{c)} $65 - x \times 3{,}2 = 45{,}8$\par
    \dongke{3}
  \end{minipage}\hfill
  \begin{minipage}[t]{0.485\linewidth}
    \textbf{b)} $x : 3{,}6 - 4{,}5 = 8{,}5$\par
    \dongke{3}
    \vspace{0.4mm}
    \textbf{d)} $(x + 15{,}5) : 4{,}5 = 8$\par
    \dongke{3}
  \end{minipage}
}

\vspace{0.8mm}

% BÀI 3: ĐỔI ĐƠN VỊ (4 Ý - 1 CỘT TỪ TRÊN XUỐNG)
\khungBaiTap{%
  \textbf{Bài 3 (1,0 điểm). Viết số thích hợp vào chỗ chấm:}\par\vspace{0.8mm}
  \noindent
  a) $15 : 50 = \dotfill\,\%$\par\vspace{1.4mm}
  \noindent
  b) $25\% \text{ của } 240\,\text{kg} = \dotfill\,\text{kg}$\par\vspace{1.4mm}
  \noindent
  c) $S_{\text{thang}} = (25 + 15) \times 10 : 2 = \dotfill\,\text{m}^2$\par\vspace{1.4mm}
  \noindent
  d) $S_{\text{tròn}} = 5 \times 5 \times 3{,}14 = \dotfill\,\text{cm}^2$\par
}

\vspace{0.8mm}

% BÀI 4: BÀI TOÁN GIẢI (11 DÒNG KẺ)
\khungBaiTap{%
  \textbf{Bài 4 (1,5 điểm).} Một mảnh đất hình thang có đáy lớn $28$ m, đáy bé $20$ m, chiều cao $15$ m. Người ta dùng $35\%$ diện tích đất để làm nhà, phần còn lại trồng cây. Tính diện tích đất trồng cây và tỉ số phần trăm diện tích đất trồng cây so với diện tích mảnh đất.
  \par\vspace{0.3mm}%
  \noindent\textit{Bài giải:}
  \dongke{11}
}

\vspace{0.8mm}

% BÀI 5: TOÁN NÂNG CAO (8 DÒNG KẺ)
\khungBaiTap{%
  \textbf{Bài 5 (1,0 điểm).} Tính nhanh:\par\vspace{0.6mm}\centerline{$K = \dfrac{0{,}25 \times 8 + 3{,}6}{1{,}2 + 0{,}8 \times 2}$}
  \par\vspace{0.3mm}%
  \noindent\textit{Bài làm:}
  \dongke{8}
}

\newpage

% =============================================================================
% ========================= MÃ ĐỀ: 502 ========================================
% =============================================================================

% -------------------- TRANG 1: MẶT TRƯỚC --------------------
\noindent
\begin{tabularx}{\linewidth}{|X|c|}
\hline
\textbf{PHIẾU BÀI TẬP CUỐI TUẦN -- TOÁN LỚP 5 (TUẦN 16)} & \textbf{MÃ ĐỀ: 502} \\
\textit{Trung tâm Toán Thầy Lĩnh -- Năm học: 2026 - 2027} & \textit{Thời gian: 45 phút} \\
\hline
\end{tabularx}

\vspace{1mm}

\noindent
\begin{tabularx}{\linewidth}{|X|p{3.5cm}|}
\hline
\rule[-2mm]{0pt}{7.5mm}\textbf{Họ và tên:} \dotfill & \textbf{Lớp:} \makebox[2cm]{\dotfill} \\
\hline
\end{tabularx}

\vspace{1mm}

\noindent
\begin{tabularx}{\linewidth}{|p{3cm}|X|}
\hline
\centering\textbf{ĐIỂM SỐ} & \centering\textbf{LỜI NHẬN XÉT CỦA THẦY CÔ} \tabularnewline
\hline
\rule{0pt}{1.1cm} & \tabularnewline
\hline
\end{tabularx}

\vspace{1mm}

\noindent\textbf{BẢNG ĐIỀN ĐÁP ÁN TRẮC NGHIỆM (Mỗi câu đúng 0,25 điểm):}
\vspace{0.6mm}

\begin{tabularx}{\linewidth}{|c|*{12}{>{\centering\arraybackslash}X|}}
\hline
\textbf{Câu} & \textbf{1} & \textbf{2} & \textbf{3} & \textbf{4} & \textbf{5} & \textbf{6} & \textbf{7} & \textbf{8} & \textbf{9} & \textbf{10} & \textbf{11} & \textbf{12} \\
\hline
\textbf{Đáp án} & & & & & & & & & & & & \\
\hline
\end{tabularx}

\vspace{1.2mm}

\noindent\textbf{\large PHẦN I. TRẮC NGHIỆM KHÁCH QUAN (3,0 điểm)}

\cauTN{1}{Kết quả của phép tính $32{,}5 \times 0{,}4$ là:}{$13$}{$130$}{$1{,}3$}{$13{,}5$}
\cauTN{2}{Giá trị của biểu thức $18{,}4 + 5{,}6 \times 3$ là:}{$72$}{$35{,}2$}{$36{,}2$}{$34{,}2$}
\cauTN{3}{Tỉ số của hai số $4$ và $7$ viết dưới dạng phân số là:}{$\dfrac{7}{4}$}{$4{,}7$}{$\dfrac{4}{7}$}{$0{,}47$}
\cauTN{4}{Phân số $\dfrac{4}{5}$ viết dưới dạng tỉ số phần trăm là:}{$80\%$}{$75\%$}{$85\%$}{$70\%$}
\cauTN{5}{Tỉ số phần trăm của $18$ và $72$ là:}{$20\%$}{$25\%$}{$30\%$}{$15\%$}
\cauTN{6}{$45\%$ của $160\,\text{kg}$ là:}{$72\,\text{kg}$}{$70\,\text{kg}$}{$75\,\text{kg}$}{$68\,\text{kg}$}
\cauTN{7}{Biết $30\%$ của một số là $45$. Số đó là:}{$15$}{$135$}{$150$}{$120$}
\cauTN{8}{Một chiếc xe đạp giá $1\,500\,000$ đồng, được giảm giá $10\%$. Số tiền được giảm là:}{$150\,000$ đồng}{$100\,000$ đồng}{$200\,000$ đồng}{$120\,000$ đồng}
\cauTN{9}{Một hình thang có hai đáy $22\,\text{cm}$ và $14\,\text{cm}$, chiều cao $10\,\text{cm}$. Diện tích hình thang là:}{$360\,\text{cm}^2$}{$180\,\text{cm}^2$}{$190\,\text{cm}^2$}{$200\,\text{cm}^2$}
\cauTN{10}{Một hình tròn có đường kính $12\,\text{cm}$. Diện tích hình tròn đó là:}{$452{,}16\,\text{cm}^2$}{$226{,}08\,\text{cm}^2$}{$113{,}04\,\text{cm}^2$}{$37{,}68\,\text{cm}^2$}
\cauTNhaiDong{11}{Mảnh đất $600\,\text{m}^2$, dành $35\%$ diện tích làm nhà, phần còn lại trồng hoa. Diện tích đất trồng hoa là:}{$390\,\text{m}^2$}{$210\,\text{m}^2$}{$350\,\text{m}^2$}{$400\,\text{m}^2$}
\cauTNhaiDong{12}{Lớp 5B có $35$ học sinh, trong đó có $21$ bạn đạt điểm tốt. Tỉ số phần trăm của bạn đạt điểm tốt là:}{$55\%$}{$60\%$}{$65\%$}{$50\%$}

\vspace{0.4mm}
\noindent\textbf{\large PHẦN II. TỰ LUẬN (7,0 điểm)}

\noindent\textbf{Bài 1 (2,0 điểm). Đặt tính rồi tính:}
\vspace{0.4mm}

\noindent
\begin{tabularx}{\linewidth}{|X|X|X|X|}
\hline
\oDatTinh[4]{a)}{$56{,}48 + 42{,}7$} & 
\oDatTinh[4]{b)}{$104{,}5 - 58{,}75$} & 
\oDatTinh[4]{c)}{$38{,}5 \times 3{,}4$} & 
\oDatTinh[4]{d)}{$136{,}8 : 3{,}6$} \tabularnewline
\hline
\end{tabularx}

\newpage
% -------------------- TRANG 2: MẶT SAU --------------------
\noindent
\begin{tabularx}{\linewidth}{|X|c|}
\hline
\textbf{PHẦN TỰ LUẬN (TIẾP THEO) -- PHIẾU BÀI TẬP TUẦN 16} & \textbf{MÃ ĐỀ: 502} \\
\hline
\end{tabularx}

\vspace{0.8mm}

% BÀI 2: TÌM X (4 Ý)
\khungBaiTap{%
  \textbf{Bài 2 (1,5 điểm). Tìm $x$:}\par\vspace{0.6mm}
  \noindent
  \begin{minipage}[t]{0.485\linewidth}
    \textbf{a)} $x \times 5{,}2 + 15{,}6 = 57{,}2$\par
    \dongke{3}
    \vspace{0.4mm}
    \textbf{c)} $78 - x \times 4{,}5 = 46{,}5$\par
    \dongke{3}
  \end{minipage}\hfill
  \begin{minipage}[t]{0.485\linewidth}
    \textbf{b)} $x : 4{,}5 - 6{,}2 = 9{,}8$\par
    \dongke{3}
    \vspace{0.4mm}
    \textbf{d)} $(x + 18{,}5) : 5{,}2 = 7$\par
    \dongke{3}
  \end{minipage}
}

\vspace{0.8mm}

% BÀI 3: ĐỔI ĐƠN VỊ (4 Ý - 1 CỘT TỪ TRÊN XUỐNG)
\khungBaiTap{%
  \textbf{Bài 3 (1,0 điểm). Viết số thích hợp vào chỗ chấm:}\par\vspace{0.8mm}
  \noindent
  a) $24 : 80 = \dotfill\,\%$\par\vspace{1.4mm}
  \noindent
  b) $35\% \text{ của } 200\,\text{kg} = \dotfill\,\text{kg}$\par\vspace{1.4mm}
  \noindent
  c) $S_{\text{thang}} = (28 + 16) \times 12 : 2 = \dotfill\,\text{m}^2$\par\vspace{1.4mm}
  \noindent
  d) $S_{\text{tròn}} = 6 \times 6 \times 3{,}14 = \dotfill\,\text{cm}^2$\par
}

\vspace{0.8mm}

% BÀI 4: BÀI TOÁN GIẢI (11 DÒNG KẺ)
\khungBaiTap{%
  \textbf{Bài 4 (1,5 điểm).} Một khu vườn hình thang có đáy lớn $36$ m, đáy bé $24$ m, chiều cao $18$ m. Người ta dành $40\%$ diện tích để làm nhà, phần còn lại trồng rau. Tính diện tích trồng rau và tỉ số phần trăm diện tích trồng rau so với diện tích khu vườn.
  \par\vspace{0.3mm}%
  \noindent\textit{Bài giải:}
  \dongke{11}
}

\vspace{0.8mm}

% BÀI 5: TOÁN NÂNG CAO (8 DÒNG KẺ)
\khungBaiTap{%
  \textbf{Bài 5 (1,0 điểm).} Tính nhanh:\par\vspace{0.6mm}\centerline{$K = \dfrac{1{,}25 \times 4 + 2{,}5}{0{,}5 \times 6 + 1{,}5}$}
  \par\vspace{0.3mm}%
  \noindent\textit{Bài làm:}
  \dongke{8}
}

\newpage

% =============================================================================
% ========================= MÃ ĐỀ: 503 ========================================
% =============================================================================

% -------------------- TRANG 1: MẶT TRƯỚC --------------------
\noindent
\begin{tabularx}{\linewidth}{|X|c|}
\hline
\textbf{PHIẾU BÀI TẬP CUỐI TUẦN -- TOÁN LỚP 5 (TUẦN 16)} & \textbf{MÃ ĐỀ: 503} \\
\textit{Trung tâm Toán Thầy Lĩnh -- Năm học: 2026 - 2027} & \textit{Thời gian: 45 phút} \\
\hline
\end{tabularx}

\vspace{1mm}

\noindent
\begin{tabularx}{\linewidth}{|X|p{3.5cm}|}
\hline
\rule[-2mm]{0pt}{7.5mm}\textbf{Họ và tên:} \dotfill & \textbf{Lớp:} \makebox[2cm]{\dotfill} \\
\hline
\end{tabularx}

\vspace{1mm}

\noindent
\begin{tabularx}{\linewidth}{|p{3cm}|X|}
\hline
\centering\textbf{ĐIỂM SỐ} & \centering\textbf{LỜI NHẬN XÉT CỦA THẦY CÔ} \tabularnewline
\hline
\rule{0pt}{1.1cm} & \tabularnewline
\hline
\end{tabularx}

\vspace{1mm}

\noindent\textbf{BẢNG ĐIỀN ĐÁP ÁN TRẮC NGHIỆM (Mỗi câu đúng 0,25 điểm):}
\vspace{0.6mm}

\begin{tabularx}{\linewidth}{|c|*{12}{>{\centering\arraybackslash}X|}}
\hline
\textbf{Câu} & \textbf{1} & \textbf{2} & \textbf{3} & \textbf{4} & \textbf{5} & \textbf{6} & \textbf{7} & \textbf{8} & \textbf{9} & \textbf{10} & \textbf{11} & \textbf{12} \\
\hline
\textbf{Đáp án} & & & & & & & & & & & & \\
\hline
\end{tabularx}

\vspace{1.2mm}

\noindent\textbf{\large PHẦN I. TRẮC NGHIỆM KHÁCH QUAN (3,0 điểm)}

\cauTN{1}{Kết quả của phép tính $45{,}5 \times 0{,}4$ là:}{$18{,}2$}{$182$}{$1{,}82$}{$18{,}02$}
\cauTN{2}{Giá trị của biểu thức $24{,}5 + 6{,}5 \times 4$ là:}{$124$}{$50{,}5$}{$51{,}5$}{$49{,}5$}
\cauTN{3}{Tỉ số của hai số $5$ và $8$ viết dưới dạng phân số là:}{$\dfrac{8}{5}$}{$5{,}8$}{$\dfrac{5}{8}$}{$0{,}58$}
\cauTN{4}{Phân số $\dfrac{1}{4}$ viết dưới dạng tỉ số phần trăm là:}{$25\%$}{$20\%$}{$30\%$}{$15\%$}
\cauTN{5}{Tỉ số phần trăm của $24$ và $80$ là:}{$25\%$}{$30\%$}{$35\%$}{$20\%$}
\cauTN{6}{$55\%$ của $200\,\text{kg}$ là:}{$110\,\text{kg}$}{$100\,\text{kg}$}{$115\,\text{kg}$}{$105\,\text{kg}$}
\cauTN{7}{Biết $40\%$ của một số là $60$. Số đó là:}{$24$}{$120$}{$150$}{$100$}
\cauTN{8}{Gửi tiết kiệm $10\,000\,000$ đồng với lãi suất $0{,}5\%$/tháng. Tiền lãi một tháng là:}{$50\,000$ đồng}{$500\,000$ đồng}{$5\,000$ đồng}{$20\,000$ đồng}
\cauTN{9}{Một hình thang có hai đáy $24\,\text{cm}$ và $16\,\text{cm}$, chiều cao $12\,\text{cm}$. Diện tích hình thang là:}{$480\,\text{cm}^2$}{$240\,\text{cm}^2$}{$250\,\text{cm}^2$}{$260\,\text{cm}^2$}
\cauTN{10}{Một hình tròn có đường kính $14\,\text{cm}$. Diện tích hình tròn đó là:}{$615{,}44\,\text{cm}^2$}{$307{,}72\,\text{cm}^2$}{$153{,}86\,\text{cm}^2$}{$43{,}96\,\text{cm}^2$}
\cauTNhaiDong{11}{Mảnh đất $750\,\text{m}^2$, dành $30\%$ diện tích làm nhà, phần còn lại đào ao. Diện tích ao là:}{$525\,\text{m}^2$}{$225\,\text{m}^2$}{$500\,\text{m}^2$}{$550\,\text{m}^2$}
\cauTNhaiDong{12}{Một đàn gà có $80$ con, trong đó có $52$ con gà mái. Tỉ số phần trăm của gà mái trong đàn là:}{$60\%$}{$65\%$}{$70\%$}{$55\%$}

\vspace{0.4mm}
\noindent\textbf{\large PHẦN II. TỰ LUẬN (7,0 điểm)}

\noindent\textbf{Bài 1 (2,0 điểm). Đặt tính rồi tính:}
\vspace{0.4mm}

\noindent
\begin{tabularx}{\linewidth}{|X|X|X|X|}
\hline
\oDatTinh[4]{a)}{$65{,}35 + 48{,}9$} & 
\oDatTinh[4]{b)}{$118{,}4 - 69{,}75$} & 
\oDatTinh[4]{c)}{$42{,}5 \times 2{,}8$} & 
\oDatTinh[4]{d)}{$159{,}8 : 4{,}7$} \tabularnewline
\hline
\end{tabularx}

\newpage
% -------------------- TRANG 2: MẶT SAU --------------------
\noindent
\begin{tabularx}{\linewidth}{|X|c|}
\hline
\textbf{PHẦN TỰ LUẬN (TIẾP THEO) -- PHIẾU BÀI TẬP TUẦN 16} & \textbf{MÃ ĐỀ: 503} \\
\hline
\end{tabularx}

\vspace{0.8mm}

% BÀI 2: TÌM X (4 Ý)
\khungBaiTap{%
  \textbf{Bài 2 (1,5 điểm). Tìm $x$:}\par\vspace{0.6mm}
  \noindent
  \begin{minipage}[t]{0.485\linewidth}
    \textbf{a)} $x \times 6{,}4 + 18{,}2 = 69{,}4$\par
    \dongke{3}
    \vspace{0.4mm}
    \textbf{c)} $86 - x \times 5{,}4 = 48{,}2$\par
    \dongke{3}
  \end{minipage}\hfill
  \begin{minipage}[t]{0.485\linewidth}
    \textbf{b)} $x : 5{,}2 - 7{,}5 = 10{,}5$\par
    \dongke{3}
    \vspace{0.4mm}
    \textbf{d)} $(x + 22{,}5) : 6{,}4 = 6$\par
    \dongke{3}
  \end{minipage}
}

\vspace{0.8mm}

% BÀI 3: ĐỔI ĐƠN VỊ (4 Ý - 1 CỘT TỪ TRÊN XUỐNG)
\khungBaiTap{%
  \textbf{Bài 3 (1,0 điểm). Viết số thích hợp vào chỗ chấm:}\par\vspace{0.8mm}
  \noindent
  a) $32 : 80 = \dotfill\,\%$\par\vspace{1.4mm}
  \noindent
  b) $45\% \text{ của } 300\,\text{kg} = \dotfill\,\text{kg}$\par\vspace{1.4mm}
  \noindent
  c) $S_{\text{thang}} = (30 + 20) \times 14 : 2 = \dotfill\,\text{m}^2$\par\vspace{1.4mm}
  \noindent
  d) $S_{\text{tròn}} = 7 \times 7 \times 3{,}14 = \dotfill\,\text{cm}^2$\par
}

\vspace{0.8mm}

% BÀI 4: BÀI TOÁN GIẢI (11 DÒNG KẺ)
\khungBaiTap{%
  \textbf{Bài 4 (1,5 điểm).} Một thửa ruộng hình thang có đáy lớn $42$ m, đáy bé $30$ m, chiều cao $20$ m. Người ta dùng $45\%$ diện tích để làm nhà, phần còn lại trồng lúa. Tính diện tích trồng lúa và tỉ số phần trăm diện tích trồng lúa so với diện tích thửa ruộng.
  \par\vspace{0.3mm}%
  \noindent\textit{Bài giải:}
  \dongke{11}
}

\vspace{0.8mm}

% BÀI 5: TOÁN NÂNG CAO (8 DÒNG KẺ)
\khungBaiTap{%
  \textbf{Bài 5 (1,0 điểm).} Tính nhanh:\par\vspace{0.6mm}\centerline{$K = \dfrac{2{,}4 \times 5 - 4}{0{,}75 \times 8 + 0{,}5}$}
  \par\vspace{0.3mm}%
  \noindent\textit{Bài làm:}
  \dongke{8}
}

\newpage

% =============================================================================
% ========================= MÃ ĐỀ: 504 ========================================
% =============================================================================

% -------------------- TRANG 1: MẶT TRƯỚC --------------------
\noindent
\begin{tabularx}{\linewidth}{|X|c|}
\hline
\textbf{PHIẾU BÀI TẬP CUỐI TUẦN -- TOÁN LỚP 5 (TUẦN 16)} & \textbf{MÃ ĐỀ: 504} \\
\textit{Trung tâm Toán Thầy Lĩnh -- Năm học: 2026 - 2027} & \textit{Thời gian: 45 phút} \\
\hline
\end{tabularx}

\vspace{1mm}

\noindent
\begin{tabularx}{\linewidth}{|X|p{3.5cm}|}
\hline
\rule[-2mm]{0pt}{7.5mm}\textbf{Họ và tên:} \dotfill & \textbf{Lớp:} \makebox[2cm]{\dotfill} \\
\hline
\end{tabularx}

\vspace{1mm}

\noindent
\begin{tabularx}{\linewidth}{|p{3cm}|X|}
\hline
\centering\textbf{ĐIỂM SỐ} & \centering\textbf{LỜI NHẬN XÉT CỦA THẦY CÔ} \tabularnewline
\hline
\rule{0pt}{1.1cm} & \tabularnewline
\hline
\end{tabularx}

\vspace{1mm}

\noindent\textbf{BẢNG ĐIỀN ĐÁP ÁN TRẮC NGHIỆM (Mỗi câu đúng 0,25 điểm):}
\vspace{0.6mm}

\begin{tabularx}{\linewidth}{|c|*{12}{>{\centering\arraybackslash}X|}}
\hline
\textbf{Câu} & \textbf{1} & \textbf{2} & \textbf{3} & \textbf{4} & \textbf{5} & \textbf{6} & \textbf{7} & \textbf{8} & \textbf{9} & \textbf{10} & \textbf{11} & \textbf{12} \\
\hline
\textbf{Đáp án} & & & & & & & & & & & & \\
\hline
\end{tabularx}

\vspace{1.2mm}

\noindent\textbf{\large PHẦN I. TRẮC NGHIỆM KHÁCH QUAN (3,0 điểm)}

\cauTN{1}{Kết quả của phép tính $52{,}5 \times 0{,}4$ là:}{$21$}{$210$}{$2{,}1$}{$21{,}5$}
\cauTN{2}{Giá trị của biểu thức $32{,}8 + 7{,}2 \times 3$ là:}{$120$}{$54{,}4$}{$55{,}4$}{$53{,}4$}
\cauTN{3}{Tỉ số của hai số $6$ và $11$ viết dưới dạng phân số là:}{$\dfrac{11}{6}$}{$6{,}11$}{$\dfrac{6}{11}$}{$0{,}611$}
\cauTN{4}{Phân số $\dfrac{1}{2}$ viết dưới dạng tỉ số phần trăm là:}{$50\%$}{$20\%$}{$40\%$}{$25\%$}
\cauTN{5}{Tỉ số phần trăm của $28$ và $70$ là:}{$35\%$}{$40\%$}{$45\%$}{$30\%$}
\cauTN{6}{$65\%$ của $240\,\text{kg}$ là:}{$156\,\text{kg}$}{$150\,\text{kg}$}{$160\,\text{kg}$}{$152\,\text{kg}$}
\cauTN{7}{Biết $50\%$ của một số là $75$. Số đó là:}{$37{,}5$}{$125$}{$150$}{$100$}
\cauTN{8}{Một chiếc tivi giá $8\,000\,000$ đồng, được giảm giá $15\%$. Số tiền được giảm là:}{$1\,200\,000$ đồng}{$1\,000\,000$ đồng}{$1\,500\,000$ đồng}{$800\,000$ đồng}
\cauTN{9}{Một hình thang có hai đáy $26\,\text{cm}$ và $18\,\text{cm}$, chiều cao $14\,\text{cm}$. Diện tích hình thang là:}{$616\,\text{cm}^2$}{$308\,\text{cm}^2$}{$320\,\text{cm}^2$}{$330\,\text{cm}^2$}
\cauTN{10}{Một hình tròn có đường kính $16\,\text{cm}$. Diện tích hình tròn đó là:}{$803{,}84\,\text{cm}^2$}{$401{,}92\,\text{cm}^2$}{$200{,}96\,\text{cm}^2$}{$50{,}24\,\text{cm}^2$}
\cauTNhaiDong{11}{Mảnh đất $800\,\text{m}^2$, dành $45\%$ diện tích làm nhà, phần còn lại làm sân vườn. Diện tích sân vườn là:}{$440\,\text{m}^2$}{$360\,\text{m}^2$}{$450\,\text{m}^2$}{$420\,\text{m}^2$}
\cauTNhaiDong{12}{Một trang trại nuôi $120$ con lợn, đã tiêm phòng cho $84$ con. Tỉ số phần trăm số lợn đã tiêm là:}{$65\%$}{$70\%$}{$75\%$}{$60\%$}

\vspace{0.4mm}
\noindent\textbf{\large PHẦN II. TỰ LUẬN (7,0 điểm)}

\noindent\textbf{Bài 1 (2,0 điểm). Đặt tính rồi tính:}
\vspace{0.4mm}

\noindent
\begin{tabularx}{\linewidth}{|X|X|X|X|}
\hline
\oDatTinh[4]{a)}{$74{,}65 + 53{,}8$} & 
\oDatTinh[4]{b)}{$126{,}2 - 78{,}45$} & 
\oDatTinh[4]{c)}{$48{,}5 \times 3{,}2$} & 
\oDatTinh[4]{d)}{$187{,}2 : 5{,}2$} \tabularnewline
\hline
\end{tabularx}

\newpage
% -------------------- TRANG 2: MẶT SAU --------------------
\noindent
\begin{tabularx}{\linewidth}{|X|c|}
\hline
\textbf{PHẦN TỰ LUẬN (TIẾP THEO) -- PHIẾU BÀI TẬP TUẦN 16} & \textbf{MÃ ĐỀ: 504} \\
\hline
\end{tabularx}

\vspace{0.8mm}

% BÀI 2: TÌM X (4 Ý)
\khungBaiTap{%
  \textbf{Bài 2 (1,5 điểm). Tìm $x$:}\par\vspace{0.6mm}
  \noindent
  \begin{minipage}[t]{0.485\linewidth}
    \textbf{a)} $x \times 7{,}2 + 24{,}6 = 89{,}4$\par
    \dongke{3}
    \vspace{0.4mm}
    \textbf{c)} $95 - x \times 6{,}2 = 51{,}6$\par
    \dongke{3}
  \end{minipage}\hfill
  \begin{minipage}[t]{0.485\linewidth}
    \textbf{b)} $x : 6{,}4 - 8{,}5 = 11{,}5$\par
    \dongke{3}
    \vspace{0.4mm}
    \textbf{d)} $(x + 25{,}5) : 7{,}2 = 5$\par
    \dongke{3}
  \end{minipage}
}

\vspace{0.8mm}

% BÀI 3: ĐỔI ĐƠN VỊ (4 Ý - 1 CỘT TỪ TRÊN XUỐNG)
\khungBaiTap{%
  \textbf{Bài 3 (1,0 điểm). Viết số thích hợp vào chỗ chấm:}\par\vspace{0.8mm}
  \noindent
  a) $35 : 70 = \dotfill\,\%$\par\vspace{1.4mm}
  \noindent
  b) $60\% \text{ của } 250\,\text{kg} = \dotfill\,\text{kg}$\par\vspace{1.4mm}
  \noindent
  c) $S_{\text{thang}} = (34 + 22) \times 15 : 2 = \dotfill\,\text{m}^2$\par\vspace{1.4mm}
  \noindent
  d) $S_{\text{tròn}} = 8 \times 8 \times 3{,}14 = \dotfill\,\text{cm}^2$\par
}

\vspace{0.8mm}

% BÀI 4: BÀI TOÁN GIẢI (11 DÒNG KẺ)
\khungBaiTap{%
  \textbf{Bài 4 (1,5 điểm).} Một mảnh đất hình thang có đáy lớn $32$ m, đáy bé $18$ m, chiều cao $20$ m. Người ta dùng $25\%$ diện tích để làm nhà, phần còn lại trồng cây ăn quả. Tính diện tích trồng cây ăn quả và tỉ số phần trăm diện tích trồng cây ăn quả so với diện tích mảnh đất.
  \par\vspace{0.3mm}%
  \noindent\textit{Bài giải:}
  \dongke{11}
}

\vspace{0.8mm}

% BÀI 5: TOÁN NÂNG CAO (8 DÒNG KẺ)
\khungBaiTap{%
  \textbf{Bài 5 (1,0 điểm).} Tính nhanh:\par\vspace{0.6mm}\centerline{$K = \dfrac{1{,}5 \times 6 + 3}{0{,}4 \times 5 + 1}$}
  \par\vspace{0.3mm}%
  \noindent\textit{Bài làm:}
  \dongke{8}
}

\newpage

% =============================================================================
% ========================= MÃ ĐỀ: 505 ========================================
% =============================================================================

% -------------------- TRANG 1: MẶT TRƯỚC --------------------
\noindent
\begin{tabularx}{\linewidth}{|X|c|}
\hline
\textbf{PHIẾU BÀI TẬP CUỐI TUẦN -- TOÁN LỚP 5 (TUẦN 16)} & \textbf{MÃ ĐỀ: 505} \\
\textit{Trung tâm Toán Thầy Lĩnh -- Năm học: 2026 - 2027} & \textit{Thời gian: 45 phút} \\
\hline
\end{tabularx}

\vspace{1mm}

\noindent
\begin{tabularx}{\linewidth}{|X|p{3.5cm}|}
\hline
\rule[-2mm]{0pt}{7.5mm}\textbf{Họ và tên:} \dotfill & \textbf{Lớp:} \makebox[2cm]{\dotfill} \\
\hline
\end{tabularx}

\vspace{1mm}

\noindent
\begin{tabularx}{\linewidth}{|p{3cm}|X|}
\hline
\centering\textbf{ĐIỂM SỐ} & \centering\textbf{LỜI NHẬN XÉT CỦA THẦY CÔ} \tabularnewline
\hline
\rule{0pt}{1.1cm} & \tabularnewline
\hline
\end{tabularx}

\vspace{1mm}

\noindent\textbf{BẢNG ĐIỀN ĐÁP ÁN TRẮC NGHIỆM (Mỗi câu đúng 0,25 điểm):}
\vspace{0.6mm}

\begin{tabularx}{\linewidth}{|c|*{12}{>{\centering\arraybackslash}X|}}
\hline
\textbf{Câu} & \textbf{1} & \textbf{2} & \textbf{3} & \textbf{4} & \textbf{5} & \textbf{6} & \textbf{7} & \textbf{8} & \textbf{9} & \textbf{10} & \textbf{11} & \textbf{12} \\
\hline
\textbf{Đáp án} & & & & & & & & & & & & \\
\hline
\end{tabularx}

\vspace{1.2mm}

\noindent\textbf{\large PHẦN I. TRẮC NGHIỆM KHÁCH QUAN (3,0 điểm)}

\cauTN{1}{Kết quả của phép tính $64{,}5 \times 0{,}4$ là:}{$25{,}8$}{$258$}{$2{,}58$}{$25{,}08$}
\cauTN{2}{Giá trị của biểu thức $41{,}5 + 8{,}5 \times 3$ là:}{$150$}{$67$}{$68$}{$66$}
\cauTN{3}{Tỉ số của hai số $7$ và $12$ viết dưới dạng phân số là:}{$\dfrac{12}{7}$}{$7{,}12$}{$\dfrac{7}{12}$}{$0{,}712$}
\cauTN{4}{Phân số $\dfrac{3}{5}$ viết dưới dạng tỉ số phần trăm là:}{$60\%$}{$50\%$}{$70\%$}{$55\%$}
\cauTN{5}{Tỉ số phần trăm của $36$ và $90$ là:}{$35\%$}{$40\%$}{$45\%$}{$30\%$}
\cauTN{6}{$75\%$ của $280\,\text{kg}$ là:}{$210\,\text{kg}$}{$200\,\text{kg}$}{$220\,\text{kg}$}{$205\,\text{kg}$}
\cauTN{7}{Biết $60\%$ của một số là $90$. Số đó là:}{$54$}{$120$}{$150$}{$100$}
\cauTN{8}{Một cửa hàng nhập quạt máy giá $400\,000$ đồng, bán lãi $15\%$. Số tiền lãi là:}{$60\,000$ đồng}{$50\,000$ đồng}{$70\,000$ đồng}{$45\,000$ đồng}
\cauTN{9}{Một hình thang có hai đáy $28\,\text{cm}$ và $20\,\text{cm}$, chiều cao $15\,\text{cm}$. Diện tích hình thang là:}{$720\,\text{cm}^2$}{$360\,\text{cm}^2$}{$380\,\text{cm}^2$}{$390\,\text{cm}^2$}
\cauTN{10}{Một hình tròn có đường kính $18\,\text{cm}$. Diện tích hình tròn đó là:}{$1\\,017{,}36\,\text{cm}^2$}{$508{,}68\,\text{cm}^2$}{$254{,}34\,\text{cm}^2$}{$56{,}52\,\text{cm}^2$}
\cauTNhaiDong{11}{Mảnh đất $900\,\text{m}^2$, dành $35\%$ diện tích làm nhà, phần còn lại trồng cây ăn quả. Diện tích trồng cây là:}{$585\,\text{m}^2$}{$315\,\text{m}^2$}{$600\,\text{m}^2$}{$550\,\text{m}^2$}
\cauTNhaiDong{12}{Khối 5 có $150$ học sinh, trong đó có $105$ bạn tham gia câu lạc bộ. Tỉ số phần trăm tham gia là:}{$65\%$}{$70\%$}{$75\%$}{$60\%$}

\vspace{0.4mm}
\noindent\textbf{\large PHẦN II. TỰ LUẬN (7,0 điểm)}

\noindent\textbf{Bài 1 (2,0 điểm). Đặt tính rồi tính:}
\vspace{0.4mm}

\noindent
\begin{tabularx}{\linewidth}{|X|X|X|X|}
\hline
\oDatTinh[4]{a)}{$82{,}75 + 61{,}6$} & 
\oDatTinh[4]{b)}{$135{,}4 - 86{,}75$} & 
\oDatTinh[4]{c)}{$52{,}5 \times 2{,}8$} & 
\oDatTinh[4]{d)}{$204 : 4{,}8$} \tabularnewline
\hline
\end{tabularx}

\newpage
% -------------------- TRANG 2: MẶT SAU --------------------
\noindent
\begin{tabularx}{\linewidth}{|X|c|}
\hline
\textbf{PHẦN TỰ LUẬN (TIẾP THEO) -- PHIẾU BÀI TẬP TUẦN 16} & \textbf{MÃ ĐỀ: 505} \\
\hline
\end{tabularx}

\vspace{0.8mm}

% BÀI 2: TÌM X (4 Ý)
\khungBaiTap{%
  \textbf{Bài 2 (1,5 điểm). Tìm $x$:}\par\vspace{0.6mm}
  \noindent
  \begin{minipage}[t]{0.485\linewidth}
    \textbf{a)} $x \times 8{,}2 + 28{,}4 = 94$\par
    \dongke{3}
    \vspace{0.4mm}
    \textbf{c)} $108 - x \times 7{,}5 = 55{,}5$\par
    \dongke{3}
  \end{minipage}\hfill
  \begin{minipage}[t]{0.485\linewidth}
    \textbf{b)} $x : 7{,}5 - 9{,}2 = 12{,}8$\par
    \dongke{3}
    \vspace{0.4mm}
    \textbf{d)} $(x + 28{,}5) : 8{,}2 = 5$\par
    \dongke{3}
  \end{minipage}
}

\vspace{0.8mm}

% BÀI 3: ĐỔI ĐƠN VỊ (4 Ý - 1 CỘT TỪ TRÊN XUỐNG)
\khungBaiTap{%
  \textbf{Bài 3 (1,0 điểm). Viết số thích hợp vào chỗ chấm:}\par\vspace{0.8mm}
  \noindent
  a) $45 : 90 = \dotfill\,\%$\par\vspace{1.4mm}
  \noindent
  b) $75\% \text{ của } 320\,\text{kg} = \dotfill\,\text{kg}$\par\vspace{1.4mm}
  \noindent
  c) $S_{\text{thang}} = (36 + 24) \times 16 : 2 = \dotfill\,\text{m}^2$\par\vspace{1.4mm}
  \noindent
  d) $S_{\text{tròn}} = 9 \times 9 \times 3{,}14 = \dotfill\,\text{cm}^2$\par
}

\vspace{0.8mm}

% BÀI 4: BÀI TOÁN GIẢI (11 DÒNG KẺ)
\khungBaiTap{%
  \textbf{Bài 4 (1,5 điểm).} Một khu đất hình thang có đáy lớn $50$ m, đáy bé $34$ m, chiều cao $24$ m. Người ta dành $30\%$ diện tích để làm nhà, phần còn lại trồng hoa. Tính diện tích trồng hoa và tỉ số phần trăm diện tích trồng hoa so với diện tích khu đất.
  \par\vspace{0.3mm}%
  \noindent\textit{Bài giải:}
  \dongke{11}
}

\vspace{0.8mm}

% BÀI 5: TOÁN NÂNG CAO (8 DÒNG KẺ)
\khungBaiTap{%
  \textbf{Bài 5 (1,0 điểm).} Tính nhanh:\par\vspace{0.6mm}\centerline{$K = \dfrac{0{,}5 \times 12 + 4}{2{,}5 \times 4 - 2}$}
  \par\vspace{0.3mm}%
  \noindent\textit{Bài làm:}
  \dongke{8}
}

\newpage

% =============================================================================
% ========================= MÃ ĐỀ: 506 ========================================
% =============================================================================

% -------------------- TRANG 1: MẶT TRƯỚC --------------------
\noindent
\begin{tabularx}{\linewidth}{|X|c|}
\hline
\textbf{PHIẾU BÀI TẬP CUỐI TUẦN -- TOÁN LỚP 5 (TUẦN 16)} & \textbf{MÃ ĐỀ: 506} \\
\textit{Trung tâm Toán Thầy Lĩnh -- Năm học: 2026 - 2027} & \textit{Thời gian: 45 phút} \\
\hline
\end{tabularx}

\vspace{1mm}

\noindent
\begin{tabularx}{\linewidth}{|X|p{3.5cm}|}
\hline
\rule[-2mm]{0pt}{7.5mm}\textbf{Họ và tên:} \dotfill & \textbf{Lớp:} \makebox[2cm]{\dotfill} \\
\hline
\end{tabularx}

\vspace{1mm}

\noindent
\begin{tabularx}{\linewidth}{|p{3cm}|X|}
\hline
\centering\textbf{ĐIỂM SỐ} & \centering\textbf{LỜI NHẬN XÉT CỦA THẦY CÔ} \tabularnewline
\hline
\rule{0pt}{1.1cm} & \tabularnewline
\hline
\end{tabularx}

\vspace{1mm}

\noindent\textbf{BẢNG ĐIỀN ĐÁP ÁN TRẮC NGHIỆM (Mỗi câu đúng 0,25 điểm):}
\vspace{0.6mm}

\begin{tabularx}{\linewidth}{|c|*{12}{>{\centering\arraybackslash}X|}}
\hline
\textbf{Câu} & \textbf{1} & \textbf{2} & \textbf{3} & \textbf{4} & \textbf{5} & \textbf{6} & \textbf{7} & \textbf{8} & \textbf{9} & \textbf{10} & \textbf{11} & \textbf{12} \\
\hline
\textbf{Đáp án} & & & & & & & & & & & & \\
\hline
\end{tabularx}

\vspace{1.2mm}

\noindent\textbf{\large PHẦN I. TRẮC NGHIỆM KHÁCH QUAN (3,0 điểm)}

\cauTN{1}{Kết quả của phép tính $72{,}5 \times 0{,}4$ là:}{$29$}{$290$}{$2{,}9$}{$29{,}5$}
\cauTN{2}{Giá trị của biểu thức $48{,}6 + 9{,}4 \times 3$ là:}{$174$}{$76{,}8$}{$77{,}8$}{$75{,}8$}
\cauTN{3}{Tỉ số của hai số $8$ và $15$ viết dưới dạng phân số là:}{$\dfrac{15}{8}$}{$8{,}15$}{$\dfrac{8}{15}$}{$0{,}815$}
\cauTN{4}{Phân số $\dfrac{7}{10}$ viết dưới dạng tỉ số phần trăm là:}{$70\%$}{$60\%$}{$75\%$}{$80\%$}
\cauTN{5}{Tỉ số phần trăm của $45$ và $150$ là:}{$25\%$}{$30\%$}{$35\%$}{$20\%$}
\cauTN{6}{$80\%$ của $350\,\text{kg}$ là:}{$280\,\text{kg}$}{$270\,\text{kg}$}{$290\,\text{kg}$}{$275\,\text{kg}$}
\cauTN{7}{Biết $75\%$ của một số là $120$. Số đó là:}{$90$}{$150$}{$160$}{$140$}
\cauTN{8}{Một đôi giày giá $600\,000$ đồng, được giảm giá $20\%$. Số tiền được giảm là:}{$120\,000$ đồng}{$100\,000$ đồng}{$150\,000$ đồng}{$80\,000$ đồng}
\cauTN{9}{Một hình thang có hai đáy $32\,\text{cm}$ và $22\,\text{cm}$, chiều cao $16\,\text{cm}$. Diện tích hình thang là:}{$864\,\text{cm}^2$}{$432\,\text{cm}^2$}{$440\,\text{cm}^2$}{$450\,\text{cm}^2$}
\cauTN{10}{Một hình tròn có đường kính $20\,\text{cm}$. Diện tích hình tròn đó là:}{$1\\,256\,\text{cm}^2$}{$628\,\text{cm}^2$}{$314\,\text{cm}^2$}{$62{,}8\,\text{cm}^2$}
\cauTNhaiDong{11}{Mảnh đất $1\\,000\,\text{m}^2$, dành $30\%$ diện tích làm nhà, phần còn lại làm sân bóng mini. Diện tích sân bóng là:}{$700\,\text{m}^2$}{$300\,\text{m}^2$}{$650\,\text{m}^2$}{$750\,\text{m}^2$}
\cauTNhaiDong{12}{Một công ty có $200$ nhân viên, trong đó có $130$ nhân viên nam. Tỉ số phần trăm nhân viên nam là:}{$60\%$}{$65\%$}{$70\%$}{$55\%$}

\vspace{0.4mm}
\noindent\textbf{\large PHẦN II. TỰ LUẬN (7,0 điểm)}

\noindent\textbf{Bài 1 (2,0 điểm). Đặt tính rồi tính:}
\vspace{0.4mm}

\noindent
\begin{tabularx}{\linewidth}{|X|X|X|X|}
\hline
\oDatTinh[4]{a)}{$91{,}65 + 68{,}8$} & 
\oDatTinh[4]{b)}{$146{,}5 - 94{,}75$} & 
\oDatTinh[4]{c)}{$58{,}5 \times 3{,}2$} & 
\oDatTinh[4]{d)}{$228 : 4{,}8$} \tabularnewline
\hline
\end{tabularx}

\newpage
% -------------------- TRANG 2: MẶT SAU --------------------
\noindent
\begin{tabularx}{\linewidth}{|X|c|}
\hline
\textbf{PHẦN TỰ LUẬN (TIẾP THEO) -- PHIẾU BÀI TẬP TUẦN 16} & \textbf{MÃ ĐỀ: 506} \\
\hline
\end{tabularx}

\vspace{0.8mm}

% BÀI 2: TÌM X (4 Ý)
\khungBaiTap{%
  \textbf{Bài 2 (1,5 điểm). Tìm $x$:}\par\vspace{0.6mm}
  \noindent
  \begin{minipage}[t]{0.485\linewidth}
    \textbf{a)} $x \times 9{,}2 + 31{,}6 = 105{,}2$\par
    \dongke{3}
    \vspace{0.4mm}
    \textbf{c)} $125 - x \times 8{,}5 = 65{,}5$\par
    \dongke{3}
  \end{minipage}\hfill
  \begin{minipage}[t]{0.485\linewidth}
    \textbf{b)} $x : 8{,}4 - 10{,}5 = 14{,}5$\par
    \dongke{3}
    \vspace{0.4mm}
    \textbf{d)} $(x + 32{,}5) : 9{,}2 = 5$\par
    \dongke{3}
  \end{minipage}
}

\vspace{0.8mm}

% BÀI 3: ĐỔI ĐƠN VỊ (4 Ý - 1 CỘT TỪ TRÊN XUỐNG)
\khungBaiTap{%
  \textbf{Bài 3 (1,0 điểm). Viết số thích hợp vào chỗ chấm:}\par\vspace{0.8mm}
  \noindent
  a) $56 : 140 = \dotfill\,\%$\par\vspace{1.4mm}
  \noindent
  b) $80\% \text{ của } 400\,\text{kg} = \dotfill\,\text{kg}$\par\vspace{1.4mm}
  \noindent
  c) $S_{\text{thang}} = (40 + 26) \times 18 : 2 = \dotfill\,\text{m}^2$\par\vspace{1.4mm}
  \noindent
  d) $S_{\text{tròn}} = 10 \times 10 \times 3{,}14 = \dotfill\,\text{cm}^2$\par
}

\vspace{0.8mm}

% BÀI 4: BÀI TOÁN GIẢI (11 DÒNG KẺ)
\khungBaiTap{%
  \textbf{Bài 4 (1,5 điểm).} Một thửa đất hình thang có đáy lớn $45$ m, đáy bé $27$ m, chiều cao $16$ m. Người ta dùng $50\%$ diện tích để làm nhà, phần còn lại trồng cây cảnh. Tính diện tích trồng cây cảnh và tỉ số phần trăm diện tích trồng cây cảnh so với diện tích thửa đất.
  \par\vspace{0.3mm}%
  \noindent\textit{Bài giải:}
  \dongke{11}
}

\vspace{0.8mm}

% BÀI 5: TOÁN NÂNG CAO (8 DÒNG KẺ)
\khungBaiTap{%
  \textbf{Bài 5 (1,0 điểm).} Tính nhanh:\par\vspace{0.6mm}\centerline{$K = \dfrac{3{,}6 \times 5 + 2}{0{,}9 \times 4 + 1{,}4}$}
  \par\vspace{0.3mm}%
  \noindent\textit{Bài làm:}
  \dongke{8}
}

\newpage

% =============================================================================
% ========================= MÃ ĐỀ: 507 ========================================
% =============================================================================

% -------------------- TRANG 1: MẶT TRƯỚC --------------------
\noindent
\begin{tabularx}{\linewidth}{|X|c|}
\hline
\textbf{PHIẾU BÀI TẬP CUỐI TUẦN -- TOÁN LỚP 5 (TUẦN 16)} & \textbf{MÃ ĐỀ: 507} \\
\textit{Trung tâm Toán Thầy Lĩnh -- Năm học: 2026 - 2027} & \textit{Thời gian: 45 phút} \\
\hline
\end{tabularx}

\vspace{1mm}

\noindent
\begin{tabularx}{\linewidth}{|X|p{3.5cm}|}
\hline
\rule[-2mm]{0pt}{7.5mm}\textbf{Họ và tên:} \dotfill & \textbf{Lớp:} \makebox[2cm]{\dotfill} \\
\hline
\end{tabularx}

\vspace{1mm}

\noindent
\begin{tabularx}{\linewidth}{|p{3cm}|X|}
\hline
\centering\textbf{ĐIỂM SỐ} & \centering\textbf{LỜI NHẬN XÉT CỦA THẦY CÔ} \tabularnewline
\hline
\rule{0pt}{1.1cm} & \tabularnewline
\hline
\end{tabularx}

\vspace{1mm}

\noindent\textbf{BẢNG ĐIỀN ĐÁP ÁN TRẮC NGHIỆM (Mỗi câu đúng 0,25 điểm):}
\vspace{0.6mm}

\begin{tabularx}{\linewidth}{|c|*{12}{>{\centering\arraybackslash}X|}}
\hline
\textbf{Câu} & \textbf{1} & \textbf{2} & \textbf{3} & \textbf{4} & \textbf{5} & \textbf{6} & \textbf{7} & \textbf{8} & \textbf{9} & \textbf{10} & \textbf{11} & \textbf{12} \\
\hline
\textbf{Đáp án} & & & & & & & & & & & & \\
\hline
\end{tabularx}

\vspace{1.2mm}

\noindent\textbf{\large PHẦN I. TRẮC NGHIỆM KHÁCH QUAN (3,0 điểm)}

\cauTN{1}{Kết quả của phép tính $84{,}5 \times 0{,}4$ là:}{$33{,}8$}{$338$}{$3{,}38$}{$33{,}08$}
\cauTN{2}{Giá trị của biểu thức $54{,}8 + 11{,}2 \times 3$ là:}{$198$}{$88{,}4$}{$89{,}4$}{$87{,}4$}
\cauTN{3}{Tỉ số của hai số $9$ và $16$ viết dưới dạng phân số là:}{$\dfrac{16}{9}$}{$9{,}16$}{$\dfrac{9}{16}$}{$0{,}916$}
\cauTN{4}{Phân số $\dfrac{9}{10}$ viết dưới dạng tỉ số phần trăm là:}{$90\%$}{$80\%$}{$85\%$}{$95\%$}
\cauTN{5}{Tỉ số phần trăm của $54$ và $180$ là:}{$25\%$}{$30\%$}{$35\%$}{$20\%$}
\cauTN{6}{$85\%$ của $400\,\text{kg}$ là:}{$340\,\text{kg}$}{$330\,\text{kg}$}{$350\,\text{kg}$}{$335\,\text{kg}$}
\cauTN{7}{Biết $80\%$ của một số là $160$. Số đó là:}{$128$}{$180$}{$200$}{$190$}
\cauTN{8}{Một người bán hàng lãi $12\%$ trên vốn $5\,000\,000$ đồng. Số tiền lãi là:}{$600\,000$ đồng}{$500\,000$ đồng}{$700\,000$ đồng}{$550\,000$ đồng}
\cauTN{9}{Một hình thang có hai đáy $36\,\text{cm}$ và $24\,\text{cm}$, chiều cao $18\,\text{cm}$. Diện tích hình thang là:}{$1\\,080\,\text{cm}^2$}{$540\,\text{cm}^2$}{$550\,\text{cm}^2$}{$560\,\text{cm}^2$}
\cauTN{10}{Một hình tròn có đường kính $24\,\text{cm}$. Diện tích hình tròn đó là:}{$1\\,808{,}64\,\text{cm}^2$}{$904{,}32\,\text{cm}^2$}{$452{,}16\,\text{cm}^2$}{$75{,}36\,\text{cm}^2$}
\cauTNhaiDong{11}{Mảnh đất $1\\,200\,\text{m}^2$, dành $25\%$ diện tích làm nhà, phần còn lại trồng cây cảnh. Diện tích trồng cây cảnh là:}{$900\,\text{m}^2$}{$300\,\text{m}^2$}{$850\,\text{m}^2$}{$950\,\text{m}^2$}
\cauTNhaiDong{12}{Một cửa hàng nhập $250$ chiếc bút, đã bán được $175$ chiếc. Tỉ số phần trăm số bút đã bán là:}{$65\%$}{$70\%$}{$75\%$}{$60\%$}

\vspace{0.4mm}
\noindent\textbf{\large PHẦN II. TỰ LUẬN (7,0 điểm)}

\noindent\textbf{Bài 1 (2,0 điểm). Đặt tính rồi tính:}
\vspace{0.4mm}

\noindent
\begin{tabularx}{\linewidth}{|X|X|X|X|}
\hline
\oDatTinh[4]{a)}{$98{,}75 + 74{,}6$} & 
\oDatTinh[4]{b)}{$158{,}2 - 96{,}85$} & 
\oDatTinh[4]{c)}{$64{,}5 \times 3{,}4$} & 
\oDatTinh[4]{d)}{$255 : 5{,}1$} \tabularnewline
\hline
\end{tabularx}

\newpage
% -------------------- TRANG 2: MẶT SAU --------------------
\noindent
\begin{tabularx}{\linewidth}{|X|c|}
\hline
\textbf{PHẦN TỰ LUẬN (TIẾP THEO) -- PHIẾU BÀI TẬP TUẦN 16} & \textbf{MÃ ĐỀ: 507} \\
\hline
\end{tabularx}

\vspace{0.8mm}

% BÀI 2: TÌM X (4 Ý)
\khungBaiTap{%
  \textbf{Bài 2 (1,5 điểm). Tìm $x$:}\par\vspace{0.6mm}
  \noindent
  \begin{minipage}[t]{0.485\linewidth}
    \textbf{a)} $x \times 9{,}5 + 35{,}8 = 121{,}3$\par
    \dongke{3}
    \vspace{0.4mm}
    \textbf{c)} $140 - x \times 9{,}2 = 75{,}6$\par
    \dongke{3}
  \end{minipage}\hfill
  \begin{minipage}[t]{0.485\linewidth}
    \textbf{b)} $x : 9{,}2 - 12{,}5 = 17{,}5$\par
    \dongke{3}
    \vspace{0.4mm}
    \textbf{d)} $(x + 36{,}5) : 9{,}5 = 5$\par
    \dongke{3}
  \end{minipage}
}

\vspace{0.8mm}

% BÀI 3: ĐỔI ĐƠN VỊ (4 Ý - 1 CỘT TỪ TRÊN XUỐNG)
\khungBaiTap{%
  \textbf{Bài 3 (1,0 điểm). Viết số thích hợp vào chỗ chấm:}\par\vspace{0.8mm}
  \noindent
  a) $72 : 180 = \dotfill\,\%$\par\vspace{1.4mm}
  \noindent
  b) $85\% \text{ của } 500\,\text{kg} = \dotfill\,\text{kg}$\par\vspace{1.4mm}
  \noindent
  c) $S_{\text{thang}} = (45 + 30) \times 20 : 2 = \dotfill\,\text{m}^2$\par\vspace{1.4mm}
  \noindent
  d) $S_{\text{tròn}} = 12 \times 12 \times 3{,}14 = \dotfill\,\text{cm}^2$\par
}

\vspace{0.8mm}

% BÀI 4: BÀI TOÁN GIẢI (11 DÒNG KẺ)
\khungBaiTap{%
  \textbf{Bài 4 (1,5 điểm).} Một mảnh vườn hình thang có đáy lớn $38$ m, đáy bé $26$ m, chiều cao $15$ m. Người ta dùng $20\%$ diện tích để làm nhà, phần còn lại trồng rau sạch. Tính diện tích trồng rau sạch và tỉ số phần trăm diện tích trồng rau sạch so với diện tích mảnh vườn.
  \par\vspace{0.3mm}%
  \noindent\textit{Bài giải:}
  \dongke{11}
}

\vspace{0.8mm}

% BÀI 5: TOÁN NÂNG CAO (8 DÒNG KẺ)
\khungBaiTap{%
  \textbf{Bài 5 (1,0 điểm).} Tính nhanh:\par\vspace{0.6mm}\centerline{$K = \dfrac{2{,}5 \times 8 - 5}{1{,}25 \times 4 + 1}$}
  \par\vspace{0.3mm}%
  \noindent\textit{Bài làm:}
  \dongke{8}
}

\newpage
\end{document}
